\chapter{Guía de etiquetado P1--P4}
\label{cap:anexoc}

\section*{C.1. Finalidad del anexo}

Este anexo documenta la guía empleada para construir mediante supervisión débil
la escala académica P1--P4. El Registro 112 Castilla y León no contiene una
prioridad oficial equivalente; por ello, la guía define un criterio reproducible para
entrenar, comparar y evaluar la Capa 2 sin atribuir validez institucional a la etiqueta.

La guía se apoya en la decisión metodológica de tratar
P1--P4 como escala académica, no jurídica ni institucional.
La muestra revisada internamente por los tres autores se documenta en el
Anexo~\ref{cap:anexod}.

\section*{C.2. Principios generales}

\begin{enumerate}
    \item \textbf{Momento de primera decisión.} La etiqueta debe reflejar la urgencia
    con la información disponible al inicio, no con el desenlace final del incidente.
    \item \textbf{Prioridad como urgencia relativa.} P1--P4 ordena necesidad de
    atención temprana; no equivale a una declaración formal de situación operativa.
    \item \textbf{Precaución razonable.} Ante información incompleta, se evita tanto
    minimizar señales críticas como asumir siempre el peor caso.
    \item \textbf{Trazabilidad.} Toda etiqueta debe poder vincularse a texto,
    categoría, señales, variables o reglas.
    \item \textbf{Independencia anti-\textit{leakage}.} No se usan medios movilizados,
    pacientes atendidos, cierre ni información posterior para etiquetar la primera
    decisión.
\end{enumerate}

\section*{C.3. Definición de niveles}

\begin{table}[ht]
\centering
\caption{Definición académica de prioridades P1--P4}
\label{tab:anexoc-prioridades}
\small
\begin{tabularx}{\textwidth}{|p{1.0cm}|p{2.5cm}|X|X|}
\hline
\textbf{Nivel} & \textbf{Nombre} & \textbf{Criterio operativo} & \textbf{Ejemplos de señales} \\ \hline
P1 & Crítica & Riesgo vital confirmado o altamente probable, fallecido, atrapamiento grave, múltiples víctimas o evolución potencialmente crítica. & fallecido, inconsciente, herido grave, atrapado crítico, intoxicación grave \\ \hline
P2 & Alta & Gravedad relevante o potencial de escalado sin riesgo vital inmediato confirmado. Requiere atención rápida y seguimiento. & herido, atrapado no crítico, incendio en vivienda, rescate, varias llamadas \\ \hline
P3 & Media & Incidente que requiere gestión activa, pero sin señales claras de riesgo vital ni escalado alto. & accidente leve, incendio contenido, incidencia viaria, meteo limitada \\ \hline
P4 & Baja & Consulta, incidencia menor, asistencia diferible o caso sin urgencia operativa aparente. & consulta, información, incidencia técnica menor, sin señales críticas \\ \hline
\end{tabularx}
\end{table}

\section*{C.4. Reglas de etiquetado}

Las reglas siguientes operativizan la asignación académica como funciones de etiquetado
reproducibles. Su finalidad es producir votos débiles, no representar una verdad
institucional del servicio 112:

\begin{table}[ht]
\centering
\caption{Reglas principales para etiquetado débil}
\label{tab:anexoc-reglas}
\scriptsize
\begin{tabularx}{\textwidth}{|p{2.8cm}|p{1.2cm}|X|}
\hline
\textbf{Regla} & \textbf{Nivel mínimo} & \textbf{Criterio} \\ \hline
R-P1-FALLECIDO & P1 & El texto menciona fallecido, muerto, cadáver o equivalente. \\ \hline
R-P1-RIESGO-VITAL & P1 & Riesgo vital textual, inconsciencia, parada, herido muy grave o situación crítica. \\ \hline
R-P1-ATRAPADO-CRITICO & P1 & Atrapamiento con indicios de gravedad, rescate urgente o imposibilidad de salida. \\ \hline
R-P2-HERIDO-GRAVE & P2 & Herido grave sin confirmación suficiente para P1. \\ \hline
R-P2-INTOXICACION & P2 & Intoxicación, humo o exposición a sustancias con posible afectación personal. \\ \hline
R-P2-INCENDIO-HABITADO & P2 & Incendio en vivienda, edificio, local habitado o entorno con personas expuestas. \\ \hline
R-P2-RESCATE & P2 & Rescate o salvamento sin señal P1 confirmada. \\ \hline
R-P3-TRAFICO & P3 & Accidente de tráfico sin señales críticas ni heridos graves. \\ \hline
R-P3-INCENDIO-SIN-VICTIMAS & P3 & Incendio sin víctimas ni atrapados descritos. \\ \hline
R-P4-SIN-URGENCIA & P4 & Consulta, información o incidencia menor sin señales operativas de gravedad. \\ \hline
\end{tabularx}
\end{table}

\section*{C.5. Resolución de conflictos}

Cuando varias señales sugieren prioridades diferentes, se aplica la prioridad más alta
compatible con la información inicial. La incertidumbre no debe rebajar una prioridad:
debe reflejarse como menor confianza o como caso para revisión humana. Si la evidencia
solo permite una lectura ambigua entre dos niveles, se selecciona el nivel superior
cuando existan señales de personas afectadas, vulnerabilidad o posible escalado.

Los conflictos más relevantes son:

\begin{itemize}
    \item \textbf{P1/P2}: se decide por la presencia de riesgo vital, fallecido,
    atrapamiento crítico o gravedad personal extrema.
    \item \textbf{P2/P3}: se decide por potencial de escalado, personas expuestas,
    varias llamadas, incendio habitado o necesidad de rescate.
    \item \textbf{P3/P4}: se decide por existencia de gestión activa frente a consulta
    o incidencia menor.
\end{itemize}

\section*{C.6. Funciones de etiquetado, agregación y acuerdo}

Cuatro funciones de etiquetado heurísticas y correlacionadas forman la etiqueta final.
Todas comparten una puntuación programática de gravedad: dos incorporan la categoría,
otra la excluye y la cuarta añade patrones de vulnerabilidad, emplazamiento y negación.
Algunos nombres internos proceden del diseño experimental inicial, pero en la ejecución
evaluada no se emplean un modelo NER, un agrupamiento vectorial ni un LLM anotador. Se
trata, por tanto, de votos programáticos correlacionados y no de anotaciones humanas
independientes ni de una verdad institucional.

Los votos se agregan mediante suma ponderada. Los pesos son \(0{,}85\), \(1{,}00\),
\(0{,}90\) y \(1{,}10\), respectivamente. Si varias prioridades obtienen el mismo
peso, se selecciona la más urgente. La ejecución canónica alcanza Krippendorff
\(\alpha = 0{,}899902\); la ablación sin reglas heurísticas mantiene
\(\alpha = 0{,}834155\).

Estos valores respaldan la consistencia del procedimiento, pero no demuestran
independencia entre funciones: varias comparten patrones y señales con las variables
que posteriormente utiliza RuleFit-lite. La ablación confirma que el acuerdo no
desaparece al retirar una fuente, no que se haya eliminado por completo la circularidad.

\section*{C.7. Limitaciones}

La guía reduce subjetividad, pero no elimina la necesidad de revisión humana. La
evaluación final identifica 59 falsos negativos P1, incluidos 7 casos P1--P4, que
constituyen la prioridad de la muestra revisada. Estos casos permiten estudiar
límites de la guía, señales insuficientes o decisiones del modelo con riesgo de
infrapriorización. La muestra completa se documenta en el
Anexo~\ref{cap:anexod}.

La validación externa con operadores del 112 queda como trabajo futuro. Hasta entonces,
P1--P4 debe citarse siempre como escala académica del prototipo.
